% Generated draft: requires mathematical and source review.
\documentclass[11pt]{article}
\usepackage[margin=25mm]{geometry}
\usepackage{fontspec}
\setmainfont{Latin Modern Roman}
\usepackage{amsmath,amssymb,mathtools,unicode-math}
\setmathfont{Latin Modern Math}
\usepackage{enumitem,needspace,xurl,hyperref}
\setlength{\parindent}{0pt}
\setlength{\parskip}{5pt}
\setlength{\emergencystretch}{3em}
\newcommand{\N}{\mathbb{N}}
\newcommand{\Z}{\mathbb{Z}}
\newcommand{\Q}{\mathbb{Q}}
\newcommand{\R}{\mathbb{R}}
\newcommand{\C}{\mathbb{C}}
\newcommand{\F}{\mathbb{F}}
\newcommand{\A}{\mathbb{A}}
\newcommand{\PP}{\mathbb{P}}
\newcommand{\E}{\mathbb{E}}
\newcommand{\eps}{\varepsilon}
\newcommand{\dd}{\,\mathrm{d}}
\begin{document}
\section*{UnsolvedMath ID: 30005209}
\textbf{Source problem code:} catalogue entry
\subsection{Definitions and mathematical statement}
The supplied record does not define the nonlocal energy \(E_\gamma\), its admissible competitors, or the ambient dimension. Consequently, it does not determine a complete mathematical classification problem.

For notation only, if \(\varphi\) is a surface tension on the unit sphere \(S^{n-1}\subset\R^n\), its conventional Wulff body is
\[
W_\varphi=\bigcap_{\nu\in S^{n-1}}\{x\in\R^n:x\cdot\nu\leq\varphi(\nu)\}.
\]
Under the usual convention, a crystalline anisotropy has a polytopal Wulff body. These conventions explain the terminology; the supplied record does not specify the hypotheses on \(\varphi\), any symmetry requirements, or the normalization of \(W_\varphi\).

The requested classification has the following conditional quantifier structure. Once the intended energy, admissible class, and normalization have been supplied, let \(\mathcal A_{\mathrm{cr}}\) denote the intended class of crystalline anisotropies, let \(\mathcal C_\varphi\) denote the admissible competitors, and let \(\widehat W_\varphi\in\mathcal C_\varphi\) denote the appropriately normalized Wulff shape. Determine exactly the anisotropies \(\varphi\in\mathcal A_{\mathrm{cr}}\) for which
\[
\exists\,\gamma_0>0\quad
\forall\,\gamma\in(0,\gamma_0)\quad
\forall\,F\in\mathcal C_\varphi,
\qquad E_\gamma(\widehat W_\varphi)\leq E_\gamma(F).
\]
Here \(\gamma_0\) may depend on the anisotropy and the fixed parameters of the intended model. The notation suppresses any anisotropy dependence of the energy, as does the source. The condition concerns positive, sufficiently small \(\gamma\); it does not assert uniqueness of minimizers or a conclusion at \(\gamma=0\).

The original passage further states that the classification would require characterizing critical points of the nonlocal energy with respect to a restricted class of variations. Neither that class nor the nonlocal term is defined in the supplied record, so this qualification cannot be made more explicit here.
\subsection{Short English statement}
Classify the crystalline anisotropies whose suitably normalized Wulff shapes globally minimize \(E_\gamma\) for every sufficiently small positive \(\gamma\). The energy, competitor class, dimension, normalization, and restricted variations must be recovered from the source before this becomes a fully specified problem. The supplied record labels the problem open; this draft does not independently verify that status.
\subsection{Sources}
\begin{itemize}
\item Supplied local UnsolvedMath registry record, permanent ID \(30005209\), title: \emph{Stable Wulff Shapes for Crystalline Nonlocal Energies}; exact source problem code: \texttt{catalogue entry}. Collection: Oberwolfach Reports Open Problems. No individual proposer is supplied.
\item Supplied report citation: \emph{Calculus of Variations (2023)}, Oberwolfach Reports, pp.~2129--2193, DOI \texttt{10.4171/owr/2022/37}. \url{https://doi.org/10.4171/owr/2022/37}. The record gives proposed year 2022 and describes the report as documenting workshop open-problem discussions.
\item Supplied literature reference: Fusco--Julin--Morini, crystalline anisotropies. \url{https://doi.org/10.1016/j.na.2020.112223}. No fuller bibliographic details are supplied.
\item Supplied literature reference: original Oberwolfach report. \url{https://doi.org/10.4171/owr/2022/37}.
\item Literature assessment dated exactly \texttt{2026-08-22}: the supplied curation labels the problem open, reports results for broad symmetry classes of crystalline surface tensions, and states that the cited primary work explicitly leaves the full classification open. This is supplied dated curation derived from the OWR material, not an independently verified literature assessment; its time-sensitive status and cited claims require independent verification.
\item The registry labels the statement \texttt{corrected\_verified} and says that only formula delimiters, notation, grammar, and typography were repaired while preserving mathematical content and scope. This supplied editorial qualification does not establish mathematical completeness or correctness.
\end{itemize}
\end{document}
